\begin{summarybox}
    \textbf{\textcolor{CSummary}{一句话核心}}
    \textbf{整系数多项式的最简有理根 $p/q$：分子整除常数项，分母整除最高次项系数；验出根后用 $qx-p$ 降次。}

    \medskip
    \textbf{\textcolor{CSummary}{使用条件与公式}}
    \begin{itemize}[leftmargin=*, itemsep=0.48em, topsep=0.35em]
        \item $f(x)=a_nx^n+\cdots+a_0$，各 $a_j$ 为整数，$n\ge1$，$a_n\ne0$。
        \item 先处理 $a_0=0$ 的情况：提出 $x$，记下 $0$ 根。
        \item 对 $a_0\ne0$，若 $p/q$ 是最简有理根、$q>0$，则
            \[
                \boxed{p\mid a_0,\qquad q\mid a_n.}
            \]
    \end{itemize}

    \medskip
    \textbf{\textcolor{CSummary}{考场操作顺序}}
    \begin{itemize}[leftmargin=*, itemsep=0.45em, topsep=0.35em]
        \item \textbf{先限：} $p$ 取常数项的正负约数；$q$ 取最高次项系数的正约数；去掉重复的非最简分数。
        \item \textbf{再试：} 优先试 $1,-1$，再试其他整数和简单真分数，逐个代入 $f(x)$。
        \item \textbf{后除：} 验出 $p/q$ 才提取 $qx-p$，长除法降次后继续求根，并检查重根。
    \end{itemize}

    \medskip
    \textbf{\textcolor{CSummary}{最小反例与检查项}}
    \begin{itemize}[leftmargin=*, itemsep=0.45em, topsep=0.35em]
        \item \textbf{必要不充分：} $x^2+1$ 的候选 $\pm1$ 都不是根。
        \item \textbf{无有理根 $\ne$ 无实根：} $x^3-x-1$ 无有理根，但在 $(1,2)$ 内有实根。
        \item \textbf{无有理根 $\ne$ 不能分解：} $(x^2+1)(x^2+4)$ 有整数系数因式分解，但无有理根。
        \item \textbf{分解完成后：} 检查因式乘积、最高次项系数、常数项和总次数是否正确。
    \end{itemize}
\end{summarybox}
